% Closed P82/P84 results with observed resource accounting; not submission-ready.
\documentclass[uplatex,a4paper,twocolumn,10pt]{jsarticle}
\usepackage[margin=20mm]{geometry}
\usepackage{amsmath}
\usepackage{url}
\title{状態保持型LLMエージェントにおける\newline 情報保持と実行制御の評価}
\author{検討用別稿・著者情報未設定}
\date{}
\begin{document}
\raggedbottom
\maketitle
\noindent\textbf{あらまし案}\quad
状態保持型LLMエージェントの全24実行を解析した. 同一の公開出力が形式拒否された14生成と，構文が異なっても同じAPI要求列となった４課題37生成を確認した. 表示，実操作，最終成果を区別する必要性を示す.

\noindent\textbf{キーワード}\quad LLMエージェント，ハーネスエンジニアリング，状態保持，時間通知，費用評価

\section{問題設定と関連研究}
状態保持型LLMエージェントの費用を減らすには，操作と最終成果を損なわない制御が必要である．しかし，入力表現をそろえて出力が安定しても，有効な操作を提出できたとは限らない．CodeAct\cite{codeact}型のコード行動とAppWorld\cite{appworld}を対象に，表現の一致と実操作の対応を調べる．

時間通知と期限強制\cite{clock}，時間適応\cite{timely}，数値的非決定性\cite{nondeterminism}，ツール列・引数・出力の反復照合\cite{consistency}，ハーネス探索\cite{beyondprompts}，実行形式\cite{localprotocol}と情報提示\cite{delivery}，出典記録\cite{repro}，ID改名と正規化\cite{canonical}は既に扱われている．論理tokenと実費用\cite{costproxy}，cacheと圧縮・遅延\cite{capc,cacheeval}も区別が必要である．

本稿は，不正呼出しの文章化で通信参照が通常文へ移る経路，保存１文脈の全18出力の実行，既知４課題24実行の全426生成の記録解析を結び付ける．一致した出力の形式拒否と，異なる構文から同じ実API要求列が生じる例を示す．正規化の新方式や一般的な高速化を提案するものではない．

\section{共通の履歴処理と通知}
公開履歴と返却済み推論を$h$，呼出しIDと結果IDの対応を保って初出順に統一する処理を$C$，表示処理を$R$とする．不正引数の呼出しを$R$が文章化した後では，構造上の参照欄だけの統一は文章内IDを変更できない．共通履歴を$R(C(h))$とし，通常文，引数，コード，観測，エラー，返却推論は変えず元記録を保持する．重複・欠落・孤立参照は拒否する．長い観測の先頭ページと参照先の表示も全条件で共通とする．これは通常の正規化順序の修正であり，品質維持の保証ではない．

仕事上限$B=1800$秒と，生成直前の実残時間を整数秒にした$r$を用いる．Mは毎回$r$を表示，Hは非表示を通知，Lは$2r>B$ならH，それ以外はMとする．整数化により半分時点より１秒未満早く切り替わり得る．現在の通知だけを要求へ付け，履歴に蓄積せず，過去内容や進行中の生成は変更しない．実期限は全条件で単調時計から強制する．残token・API・生成回数の表示は共通である．

全仕事の論理費用$T=\sum_j(I_j+O_j)$は全入力，出力，制御，失敗，回復，cache入力を含む．cache除外入力，出力，API数，実時間は別に集計し，終端使用量の不明を０や別反復で補わない．

\section{比較の手順}
\subsection{通知比較と保存文脈の実行}
写真移動，電気代精算，返金，曲情報書出しの既知４課題をM/H/Lで２回ずつ初期化して実行する．全24条件で履歴処理は共通である．Mac Studio M3 Ultra・256 GiB，AppWorld 0.1.3.post1，MiniMax-M3 IQ4\_XS，llama.cpp b11669を用いる．入力枠65,536，１スロット，持続Python，型付きAPI返却，通常例外とし，温度1.0，top-p 0.95，環境seed42，モデルseed57から生成ごとに１増やす．反復１の課題別順序はMHL/HLM/LMH/LHM，反復２は各組の逆順とする．

各条件80生成，450,000論理token，1,800作業秒，写真160・他80 APIを上限とする．出力上限は4,096と残予算の小さい方，通信待ちは600秒と残時間の小さい方である．初回はcache再利用０を確認し，以後は通常のprefix cacheを使う．最終状態を独立評価器で各１回判定し，基準・結果をモデルへ返さない．Lの固定採用基準は，各反復でM/Hそれぞれに対し全件評価・使用量既知・監査一致，成功喪失０，共通成功１課題以上，全４課題の論理費用減を満たすこととする．全４比較で成立を要求し，閾値探索や成功部分だけの選択はしない．

以前のM/H反復を見た後，書出しで文章化ID差が初めて現れた保存文脈を選んだ．元入力A/B，表示前統一C/D，Aの構造化参照だけを改名したN，診断文を足したXを固定する．token列はC/DとA/Nで一致し，A/BとA/Xで異なる．各３回計18生成をseed70，温度1，出力上限4,096で固定順に測る．全要求をcold入力とし，元のwarm状態は再現しない．空出力，途中切れ，揺れ，欠測を除いて安定と判定しない．

全18出力は新規環境で各１回，先行12コード・28 APIの実要求・返答・標準出力・例外・schemaを保存記録と照合後，通常実行口へ渡す．保存返答による代用や本文からのコード救済をしない．不一致・不明時は停止し補充しない．双方の提出と実行がなければ，空API列同士を有効操作の一致としない．元のDB・プロセス全体の復元ではない．

\subsection{閉じた全24条件の記録解析}
上の24実行の全生成を事後解析する．各課題６実行の全15組，計60組を生成序数で対応付け，片側終了後も残す．公開出力は本文・関数名・引数・終了理由，構文は位置情報だけを外したAST，実API列は送信receiptを持つ全要求の名前・引数を比較する．認証値も保持し，JSON objectのキー順・整形だけを統一する．全Python内API記録を外側の実行記録へ一度だけ対応付け，元監査の費用・欠測を照合する．生成前のAPI要求・返答・例外・schema列が同じ範囲も別集計するが，Python変数やDB全体の一致とはしない．両方空のAPI列は非空要求の一致へ数えず，記録を読み直すだけでモデル・API・コード・評価を再実行しない．

\section{全仕事の結果}
\begin{table}[tb]
\centering\footnotesize\setlength{\tabcolsep}{2pt}
\caption{全24条件の成否と論理入出力token．不明生成を０へ置換せず，未評価・未実行も残す．}
\begin{tabular}{lrrr}
課題・反復 & M & H & L \\ \hline
写真移動1 & 成功 128,796 & 成功 135,628 & 成功 135,628 \\
写真移動2 & 未達 152,155 & 成功 135,628 & 成功 135,628 \\
精算1 & 成功 241,963 & 成功 211,110 & 成功 211,110 \\
精算2 & 成功 275,944 & 成功 211,110 & 成功 211,110 \\
返金1 & 成功 138,701 & 成功 102,087 & 成功 102,087 \\
返金2 & 成功 109,383 & 成功 102,087 & 成功 102,087 \\
曲情報の書出し1 & 未達 257,266 & \shortstack[r]{未達 358,638\\＋不明1} & \shortstack[r]{未達 348,419\\＋不明1} \\
曲情報の書出し2 & 未達 90,721 & \shortstack[r]{未達 358,638\\＋不明1} & \shortstack[r]{未達 346,789\\＋不明1} \\
\hline\end{tabular}

\end{table}
採用基準を満たした比較は0/4であり，両対照を両反復で上回るという固定基準は未達だった．条件を選び直した改善や，新手法の成立として扱わない．
反復1の固定基準はMに対し未達，Hに対し未達であった．
反復2の固定基準はMに対し未達，Hに対し未達であった．
反復1のM/H/Lの成功数は3/4，3/4，3/4であった．
反復2のM/H/Lの成功数は2/4，3/4，3/4であった．
全24条件で426生成，既知4,602,713論理token，使用量不明4生成，1156 API送信を記録した．所有推論器の稼働は16,263.578秒，うち読込10.061秒，評価41.470秒であった．仕事ごとの準備・作業・後片付けの合計は16,192.106秒であり，これらは重なりを持つため加算しない．
補助的な資源区間は本稿に含めず，欠測を含む全費用差を既知分だけから推定しない．

Mの反復間は全４対中4対が判定可能で，そのうち全要求一致は0対，全公開出力一致は0対であった．
Hの反復間は全４対中3対が判定可能で，そのうち全要求一致は3対，全公開出力一致は3対であった．
Lの反復間は全４対中3対が判定可能で，そのうち全要求一致は3対，全公開出力一致は3対であった．
Lで最初に実測残時間を通知した生成番号は反復1/2の順で，課題1は未到達/未到達，課題2は未到達/未到達，課題3は未到達/未到達，課題4は16/16であった．
同条件の反復間12対と，反復内の条件間24対を全て記録した．36対は実行を共有しており，独立標本数ではない．一致数は判定可能な対に限り，欠測側を一致扱いしない．

保存ログからの狭い条件付き区間は，Hの欠測２生成に正の生成進捗がなく前提不成立となった．欠測４件は残し，以下は公開レシートと非負性だけを使った事後集計である．
\begin{table}[tb]
\centering\footnotesize\setlength{\tabcolsep}{3pt}
\caption{２反復全件のLと各対照との差．下限は欠測計数の非負性だけによる．時間は実測差である．}
\begin{tabular}{lrr}
測定 & L−M & L−H \\ \hline
論理token & $\geq +197,929$ & 未確定 \\
cache除外入力 & $\geq +19,977$ & 未確定 \\
出力token & $\geq +20,204$ & 未確定 \\
作業秒 & +1,702.316 & +2.611 \\
\hline\end{tabular}
\end{table}
L−Mの論理token差は，成功同士5組では$-132,765$，失敗同士2組では$\geq +347,221$，成否不一致1組では$-16,527$であった．各区分は同じ４課題の反復を含み，独立標本ではない．成功部分の減少を全仕事へ拡張せず，Mの成功喪失も保持する．
L−Hは両側に欠測があるため既知部分の差を全費用差へ読み替えない．片側境界は信頼区間ではなく，GPU演算量，消費電力，課金，因果的な高速化も測っていない．

補助的な事後照合として，旧Hの３反復と本群Hの２反復を全４課題で比べた24対を残した．成功３課題は全要求・公開出力列が一致し，書出しの全６対は最初の13生成が一致して14生成目で差があった．本群Hの書出し両反復の14生成目は保存統一入力C/Dと要求JSONが一致したが，課題は両方未達だった．これは時期の異なる群の記録照合であり，正規化の有無を同時に割り付けた因果比較ではない．次出力診断ではcacheを無効化しており，本群のwarm状態の再現でもない．


\section{保存文脈の出力と実操作}
\begin{table}[tb]
\centering\footnotesize\setlength{\tabcolsep}{3pt}
\caption{一致数／比較可能対数．各９対は３出力ずつの組合せで，独立９課題ではない．---は双方ツール提出なし．}
\begin{tabular}{lrrrr}
入力対 & 公開出力 & コード構文 & API要求列 & 実行結果 \\ \hline
A/B & 0/9 & 0/9 & 0/9 & 0/9 \\
A/N & 9/9 & 9/9 & 9/9 & 9/9 \\
A/X & 0/9 & 0/9 & 9/9 & 0/9 \\
C/D & 9/9 & --- & --- & --- \\
\hline\end{tabular}
\end{table}
全６入力で各３出力が完全・非空で安定し，全18生成は既知422,475論理token，使用量不明０，推論器稼働4,386.052秒だった．全15入力間比較と返却messageの不透明hashを保持し，内部推論の意味内容は読まない．
全18条件で先行履歴の照合が成立した．A/N/Xの次操作はログインと曲・アルバム・プレイリスト一覧の各４ API，Bはログインと曲一覧の各２ APIだった．C/Dの６出力は本文だけでツール提出がなく，通常の実行口は形式エラーを返した．
A/Bの構文差は実API要求の差を伴った．一方，A/Xは構文が違っても引数を含むAPI要求列は同じで，標準出力は各521文字と488文字だった．構文差はAPI差と同義でなく，本文の一致も有効な操作の一致とは限らない．
実行診断は再構成504＋次42＝546 API，228コード，64.561853秒を要した．追加生成・評価は０，全所有資源を終了した．この再構成費を除いて削減率を作らず，局所操作の比較を課題成功へ広げない．


\section{全生成における表示と実要求}
全426生成では，提出した全ツールの実行記録あり397，提出なし21，構造化提出後の未完了４，出力不明４だった．397にはPython通常例外を含み，課題成功数ではない．
全1,132対応枠は両側記録998，片側終了後134であった．公開出力は992比較中441一致，構文は932中441一致，非空API列は861中446一致だった．分母の違う一致率を指標の優劣にしない．
\begin{table}[tb]
\centering\footnotesize\setlength{\tabcolsep}{3pt}
\caption{構文が異なるが非空API列が同じ比較位置．右列は生成前のAPI履歴も同じ範囲．組数は独立標本ではない．}
\begin{tabular}{lrr}
課題 & 全対応 & 前史一致 \\ \hline
写真移動 & 21 & 16 \\
電気代精算 & 1 & 0 \\
返金 & 10 & 2 \\
曲情報書出し & 1 & 1 \\
全体 & 33 & 19 \\
\hline\end{tabular}
\end{table}
同じ公開出力でも形式拒否された比較位置は19組で，全てCSVの４実行・14生成に対応した．10生成は提出なし，４生成は引数JSONが不正だった．全19組で生成前のAPI履歴も同じだった．
異なる構文・同じ非空API列の33組は，４課題16実行37生成に対応した．前史も同じ19組は３課題９実行26生成だった．37生成中30生成はPython正常終了，７生成は通常例外であり，同じ要求列を送信したことと正常完了を区別する．API履歴の一致には初回の空prefixも含む．全記録の解析は2.940039秒，追加生成・API・コード実行・評価０だった．


\section{考察と適用範囲}
出力一致には実行口で拒否された生成が含まれ，構文差だけでは実API要求の差を判定できなかった．公開出力や構文の比較には，提出・実行状態と実要求を添える必要がある．API列の一致自体も正常完了や正解を保証しない．既存研究が区別する行動再現性指標を新方式とせず，本実装での対応と失敗の具体的なデータとして位置付ける．

独立した課題は調整に用いた既知４つで，24実行，60組，生成数を独立課題数へ読み替えない．保存文脈は事後に選んだ１つで，全群解析もその観察後の二次解析である．入力・API履歴の一致は完全な実行状態の復元でなく，同じseedにも数値実行，cache，ホスト負荷の変動が残る．通知文の長さと情報の効果は分離していない．非表示は生成コードの時刻参照を禁じず，表示開始に至らないLも残す．参照統一は全条件共通なので，その因果的な品質効果は識別できない．局所API差を全仕事の削減にせず，未達・欠測と再構成費を残す．全条件と出典を公開する\footnote{\url{https://github.com/sophiria1225/AI-}}．

\section{結論}
状態保持型LLMエージェントの表現と実操作を，保存文脈の実行と全24実行の記録から評価した．公開出力が同じでも形式拒否された14生成と，構文が異なっても実API列が同じだった４課題37生成を確認した．通知比較の全４採用基準は未達だった．出力や構文の一致だけを有効操作や費用改善の代理にせず，実行状態，最終成果，失敗を含む総費用を分けて報告する必要がある．他課題への広がりと一般的な品質・費用改善，新方式の成立は未検証である．


% The verified bibliography is retained from the closed prior-review draft.
\begin{thebibliography}{99}
\raggedright
\bibitem{codeact} X. Wang, Y. Chen, L. Yuan, Y. Zhang, Y. Li, H. Peng, and H. Ji, ``Executable code actions elicit better LLM agents,'' Proc. ICML, PMLR, vol.235, pp.50208--50232, Vienna, Austria, July 2024.
\bibitem{appworld} H. Trivedi, T. Khot, M. Hartmann, R. Manku, V. Dong, E. Li, S. Gupta, A. Sabharwal, and N. Balasubramanian, ``AppWorld: a controllable world of apps and people for benchmarking interactive coding agents,'' Proc. ACL, vol.1, pp.16022--16076, Bangkok, Thailand, Aug. 2024.
\bibitem{clock} A. Wang, N. Madan, V. Sobal, M. Trager, M. Kleinman, E. Mansimov, W. Xia, and S. Soatto, ``On the Clock: Towards Punctual and Productive Time-Budgeted AI Agents,'' arXiv:2610.10833v1, 2026.
\bibitem{timely} Y. Ma, L. Li, Y. Chen, P. Li, X. Li, Q. Guo, D. Lin, and K. Chen, ``Timely Machine: Awareness of Time Makes Test-Time Scaling Agentic,'' Proc. ACL, vol.1, pp.4619--4636, San Diego, USA, July 2026.
\bibitem{nondeterminism} J. Yuan, H. Li, X. Ding, W. Xie, Y.-J. Li, W. Zhao, K. Wan, J. Shi, X. Hu, and Z. Liu, ``Understanding and Mitigating Numerical Sources of Nondeterminism in LLM Inference,'' Advances in Neural Information Processing Systems, vol.38, 2025, doi:10.52202/085713-5653.
\bibitem{repro} P. Pattnayak and A. Bhatia, ``ReproEvalCard: A Reporting Standard for Reproducible Evaluation of LLM Pipelines,'' Proc. ACL, vol.2, pp.238--249, San Diego, USA, July 2026.
\bibitem{consistency} A. Yagubyan, ``How Consistent Are LLM Agents? Measuring Behavioral Reproducibility in Multi-Step Tool-Calling Pipelines,'' arXiv:2605.28840v1, 2026.
\bibitem{beyondprompts} C. Zhao, H. Ruan, W. Chen, P.-f. Tu, U. Abbasi, and J. Hesch, ``Beyond Prompts: Measuring and Optimizing LLM Tool-Agent Harnesses,'' arXiv:2609.05736v1, 2026.
\bibitem{localprotocol} R. Agarwal, ``Measuring Local SLM Agents as Systems: Protocol Effects on Quality and Resource Use,'' preprint, 2026, doi:10.5281/zenodo.21903714.
\bibitem{delivery} K. Zhang, S. Chundur, J. P. Owusu, M. J. Abdolhosseini Qomi, and M. Raissi, ``Same Simulation, Different Evidence: Where Retrieved Tool Output Costs LLM Agents Accuracy,'' arXiv:2610.26836v1, 2026.
\bibitem{canonical} A. J. Mercado and A. Lomuscio, ``Formal Verification of Agentic Systems over Operational Data,'' arXiv:2608.03609v1, 2026.
\bibitem{costproxy} S. Weinberger and A. Hozez, ``Token Reduction Is Not Cost Reduction,'' arXiv:2607.12161v5, 2026.
\bibitem{capc} Y. Song, ``Cache-Aware Prompt Compression: A Two-Tier Cost Model for LLM API Caching,'' arXiv:2607.15516v1, 2026.
\bibitem{cacheeval} E. Lumer, F. Nizar, A. Jangiti, K. Frank, A. Gulati, M. Phadate, and V. K. Subbiah, ``Don't Break the Cache: An Evaluation of Prompt Caching for Long-Horizon Agentic Tasks,'' arXiv:2601.06007v2, 2026.
\end{thebibliography}

\clearpage
\onecolumn
\noindent\textbf{Title}\quad Evaluating Information Retention and Execution Control in Stateful LLM Agents

\noindent\textbf{Authors}\quad 未設定

\noindent\textbf{Keywords}\quad LLM agents, harness engineering, state retention, time notices, cost evaluation

\noindent\textbf{Summary draft}\quad We analyzed 24 recorded runs of a stateful LLM agent. Matching public outputs included 14 rejected generations; different syntax produced identical API requests in 37 generations across four known tasks. These observations distinguish output agreement, actual execution, and task outcomes without establishing general quality or cost improvements.
\end{document}
